\chapter{Wie weit reicht die Vorhersage}

\section{Unregelmäßig ist nicht chaotisch}

\begin{merksatz}
\boxtitel{Die gereizte Herzzelle, nachgemessen}
Kapitel 6 zeigte bei $T_p = 18$ ein Muster, das sich nicht wiederholt. Ist es Chaos? Man lasse zwei Zellen mit Abstand $10^{-8}$ starten und messe, um welchen Faktor der Abstand je Reiz wächst --- im Mittel, über $1500$ Reize, in Logarithmen:

\begin{center}
{\small
\renewcommand{\arraystretch}{1.46}
\begin{tabular}{@{}P{1.2cm}P{2.4cm}P{1.8cm}P{3.4cm}@{}}
\toprule
$T_p$ & \textbf{Muster (Kap. 6)} & $\lambda$ \textbf{je Reiz} & \textbf{Urteil}\\
\midrule
$40$ & $1\!:\!1$ & $-2{,}55$ & periodisch, stark anziehend\\
$20$ & $2\!:\!1$ & $-1{,}24$ & periodisch\\
$18{,}5$ & $2\!:\!1$ & $-0{,}39$ & periodisch, schwach\\
$18{,}2$ & unregelmäßig & $-0{,}008$ & \textbf{quasiperiodisch}\\
$18{,}0$ & unregelmäßig & $-0{,}06$ & \textbf{quasiperiodisch}\\
$17{,}8$ & unregelmäßig & $-0{,}01$ & \textbf{quasiperiodisch}\\
$17{,}5$ & $3\!:\!1$ & $-0{,}34$ & periodisch\\
\bottomrule
\end{tabular}}
\end{center}

\begin{center}
\begin{minipage}{0.9\textwidth}
\textbf{In der unregelmäßigen Zone ist $\lambda \approx 0$: Benachbarte Verläufe trennen sich nicht. Das Muster ist nicht chaotisch, sondern quasiperiodisch --- der Oszillator läuft zwischen zwei Arnold-Zungen, ohne eine zu betreten. Es sieht aus wie Chaos; es ist keines.}
\end{minipage}
\end{center}

\begin{center}
\begin{minipage}{0.88\textwidth}
\kursiv{\textbf{Das Auge unterscheidet nicht zwischen \enquote{wiederholt sich nicht} und \enquote{hängt empfindlich vom Anfang ab}.} Die Zahl in der dritten Spalte tut es: Sie ist negativ für Periodisches, null für Quasiperiodisches, positiv für Chaos. Dieses Kapitel definiert sie, zeigt, wie man sie aus Gleichungen und aus Daten gewinnt, und rechnet aus, was sie für die Vorhersage bedeutet. Für das Modell der Herzzelle mit Stößen ist das Ergebnis über alle geprüften Reizstärken dasselbe: $\lambda \le 0$. Das Chaos, das \Person{Guevara} und \Person{Glass} 1981 gemessen haben, braucht eine Phasenantwortkurve, die dieses Modell nicht hat --- eine Grenze des Modells, keine der Theorie.}
\end{minipage}
\end{center}
\end{merksatz}

\section{Der Exponent}

\begin{definition}
\boxtitelm{Definition 9.1 --- Ljapunow-Exponent}
Für einen Verlauf $x(t)$ und eine infinitesimale Störung $\delta(t)$, die sich nach der linearisierten Gleichung $\dot\delta = J(x(t))\,\delta$ entwickelt, ist der \Begriff{größte Ljapunow-Exponent}
\[ \lambda_1 = \lim_{t\to\infty}\frac1t\ln\frac{|\delta(t)|}{|\delta(0)|} . \]
Für ein System in $n$ Dimensionen gibt es $n$ Exponenten $\lambda_1 \ge \cdots \ge \lambda_n$, die Wachstumsraten in $n$ unabhängigen Richtungen; für eine Abbildung ersetzt man $t$ durch die Schrittzahl.
\end{definition}

\begin{merksatz}
\boxtitel{Satz 9.2 --- was die Exponenten sagen}
\begin{center}
{\small
\renewcommand{\arraystretch}{1.5}
\begin{tabular}{@{}P{3.4cm}P{5.4cm}@{}}
\toprule
\textbf{Stabiler Fixpunkt} & alle $\lambda_i < 0$\\
\textbf{Stabiler Zyklus} & $\lambda_1 = 0$ (entlang des Flusses), Rest $< 0$\\
\textbf{Quasiperiodisch} & $\lambda_1 = \lambda_2 = 0$, Rest $< 0$\\
\textbf{Chaos} & $\lambda_1 > 0$; für einen Fluss zudem ein $\lambda_i = 0$\\
\textbf{Summe} & $\sum\lambda_i = $ mittlere Divergenz $\overline{\nabla\cdot F}$ (Satz 1.3)\\
\bottomrule
\end{tabular}}
\end{center}

\schritt{Warum ein Exponent null ist}
Die Störung in Flussrichtung, $\delta = F(x)$, erfüllt $\dot\delta = J\,F = \frac{\mathrm d}{\mathrm dt}F(x(t))$ --- sie ist selbst eine Lösung der linearisierten Gleichung, und sie wächst nicht, weil der Verlauf beschränkt ist: Ein Zeitversatz bleibt ein Zeitversatz. $\blacksquare$

\schritt{Lorenz, $r = 28$}
$\lambda_1 = 0{,}906$, $\lambda_2 = 0$, $\lambda_3 = -14{,}57$. Summe $-13{,}67 = -\sigma - 1 - b$. \checkmark\ Die Richtung, in der Nachbarn sich trennen, dehnt um $\mathrm e^{0{,}906} = 2{,}47$ je Zeiteinheit; die Richtung quer zum Blatt zieht sich um $\mathrm e^{-14{,}57}$ zusammen.

\schritt{Logistische Abbildung}
\begin{center}
{\small
\renewcommand{\arraystretch}{1.46}
\begin{tabular}{@{}P{1.0cm}P{1.0cm}P{1.0cm}P{1.0cm}P{1.0cm}P{1.2cm}P{1.0cm}P{1.0cm}@{}}
\toprule
$r$ & $3{,}2$ & $3{,}5$ & $3{,}57$ & $3{,}7$ & $3{,}8284$ & $3{,}9$ & $4$\\
\midrule
$\lambda$ & $-0{,}92$ & $-0{,}87$ & $0{,}01$ & $0{,}35$ & $0{,}13$ & $0{,}49$ & $0{,}693$\\
\bottomrule
\end{tabular}}
\end{center}
Negativ in der Kaskade, null bei $r_\infty$, positiv dahinter, eingebrochen im Dreierfenster, bei $r = 4$ genau $\ln2$: Dort ist die Abbildung zu $x\mapsto 2x \bmod 1$ äquivalent, die jeden Schritt eine Binärstelle vernichtet.

\begin{center}
\begin{minipage}{0.88\textwidth}
\kursiv{\textbf{$\ln 2$ je Schritt heißt: ein Bit je Schritt.} Wer $x_0$ auf $50$ Bit kennt --- doppelte Genauigkeit ---, kennt $x_{50}$ auf kein Bit mehr. Die Zahl ist zugleich die Rate, mit der das System Information über seinen Anfang verliert, oder --- gleichwertig --- Information aus den unsichtbaren Stellen des Anfangs an die Oberfläche holt: die \Begriff{Kolmogorow-Sinai-Entropie}, die für eindimensionale Abbildungen gleich $\lambda$ ist und allgemein gleich der Summe der positiven Exponenten (Pesin, Auslassung 10). Chaos ist, in dieser Sprache, eine Informationsquelle mit fester Bitrate.}
\end{minipage}
\end{center}
\end{merksatz}

\section{Berechnen}

\begin{merksatz}
\boxtitel{Satz 9.3 --- das Verfahren von Benettin}
Man integriere den Verlauf $x(t)$ und daneben die linearisierte Gleichung für $\delta(t)$, normiere $\delta$ alle $\tau$ Zeiteinheiten auf Länge $\delta_0$ und merke sich die Faktoren $d_k = |\delta(k\tau)|/\delta_0$ vor der Normierung. Dann
\[ \lambda_1 = \frac{1}{K\tau}\sum_{k=1}^{K}\ln d_k . \]
Für alle $n$ Exponenten führe man $n$ Störungsvektoren mit und orthonormiere sie (Gram--Schmidt) bei jeder Normierung; der $i$-te Vektor liefert $\lambda_i$.

\schritt{Warum normieren}
Ohne Normierung wächst $\delta$ über jede Rechnergenauigkeit; und ohne Orthonormierung kippen alle Vektoren in die am stärksten wachsende Richtung und liefern alle $\lambda_1$. Das Verfahren ersetzt die linearisierte Gleichung oft durch einen zweiten Verlauf im Abstand $\delta_0$ --- so sind die Tabellen dieses Bandes gerechnet, mit $\delta_0 = 10^{-8}$ und $\tau = 0{,}5$.

\schritt{Probe}
Für Lorenz liefert das Verfahren mit $K = 2000$ den Wert $0{,}907$; der Literaturwert ist $0{,}9056$. \checkmark

\begin{center}
\begin{minipage}{0.88\textwidth}
\kursiv{\textbf{Dass der Grenzwert in Definition 9.1 existiert und für fast alle Anfangspunkte derselbe ist, ist der Satz von \Person{Oseledets} von 1968} --- der multiplikative Ergodensatz, Auslassung 11: Das Produkt vieler Jacobi-Matrizen entlang eines Verlaufs hat, wie eine Summe vieler Zufallszahlen, ein Gesetz der großen Zahlen. Es ist der Grund, warum eine einzige Zahl ein ganzes Attraktorverhalten kennzeichnet, und warum die Messung an einem Verlauf für alle gilt.}
\end{minipage}
\end{center}
\end{merksatz}

\section{Der Vorhersagehorizont}

\begin{merksatz}
\boxtitel{Satz 9.4 --- Logarithmus der Genauigkeit}
Ist der Anfang auf $\delta_0$ genau bekannt und die Vorhersage brauchbar, solange der Fehler unter $\Delta$ bleibt, so reicht sie bis
\[ t_{\max} = \frac1{\lambda_1}\ln\frac{\Delta}{\delta_0} . \]
Eine Verbesserung der Genauigkeit um den Faktor $10$ verlängert den Horizont um $\ln10/\lambda_1$ --- immer um denselben Betrag, nicht um denselben Faktor.

\begin{center}
{\small
\renewcommand{\arraystretch}{1.46}
\begin{tabular}{@{}P{2.9cm}P{1.5cm}P{2.0cm}P{2.3cm}@{}}
\toprule
\textbf{System} & $\lambda_1$ & \textbf{Verdopp\-lungs\-zeit} & \textbf{$+1$ Stelle bringt}\\
\midrule
Lorenz, $r = 28$ & $0{,}906$ & $0{,}77$ & $2{,}5$ Einheiten\\
Logistisch, $r = 4$ & $0{,}693$ & $1$ Schritt & $3{,}3$ Schritte\\
Rössler & $0{,}071$ & $9{,}8$ & $32$ Einheiten\\
Wetter (synoptisch) & $\approx 0{,}35/\mathrm d$ & $\approx 2$ Tage & $\approx 6{,}6$ Tage\\
Sonnensystem (innere Planeten) & $0{,}2$ je Mio.\,a & $3{,}5$ Mio.\,a & $11$ Mio.\,a\\
\bottomrule
\end{tabular}}
\end{center}

\begin{center}
\begin{minipage}{0.9\textwidth}
\textbf{Der Horizont ist nicht die Grenze des Wissens, sondern das Verhältnis zweier Zahlen: wie genau man anfängt und wie schnell das System vergisst. Beide kann man messen; nur eine kann man ändern.}
\end{minipage}
\end{center}

\begin{center}
\begin{minipage}{0.88\textwidth}
\kursiv{\textbf{Die Zeile des Sonnensystems ist die erstaunlichste}: \Person{Laskar} fand 1989 in einer numerischen Integration über $200$ Millionen Jahre einen Ljapunow-Exponenten von etwa $1/5$ Millionen Jahren für die inneren Planeten. Die Bahn der Erde ist auf zehn Millionen Jahre hinaus vorhersagbar, auf hundert Millionen nicht --- nicht, weil die Beobachtungen schlecht wären, sondern weil jede Verbesserung um eine Stelle nur elf Millionen Jahre bringt. Newtons Uhrwerk hat einen Horizont. Für das Wetter: Die zwei Tage Verdopplungszeit sind die Zahl, die \Person{Lorenz} 1969 schätzte und die \Person{Zhang} und Mitarbeiter 2019 mit modernen Modellen bestätigten --- und aus ihr folgt, dass eine Vorhersage über zwei Wochen hinaus durch keine Messung der Welt möglich wird. Kapitel 15 sagt, was man stattdessen tut.}
\end{minipage}
\end{center}
\end{merksatz}

\section{Der Bruchfall}

\begin{warnung}
\boxtitelr{Rauschen hat einen Exponenten --- einen unendlichen}
Man nehme eine Zeitreihe und schätze $\lambda_1$ wie aus Daten üblich (Kapitel 12): Für jeden Punkt den nächsten Nachbarn suchen, den Abstand über $k$ Schritte verfolgen, den mittleren Logarithmus gegen $k$ auftragen, die Steigung lesen.

\begin{center}
{\small
\renewcommand{\arraystretch}{1.46}
\begin{tabular}{@{}P{3.4cm}P{0.8cm}P{0.8cm}P{0.8cm}P{0.8cm}P{0.8cm}P{0.8cm}@{}}
\toprule
\textbf{Steigung bei Schritt} $k$ & $1$ & $2$ & $3$ & $4$ & $6$ & $8$\\
\midrule
Logistisch, $r = 4$ & $0{,}69$ & $0{,}69$ & $0{,}69$ & $0{,}69$ & $0{,}70$ & $0{,}67$\\
Rauschen, AR(1), $0{,}9$ & $1{,}54$ & $0{,}64$ & $0{,}36$ & $0{,}15$ & $0{,}06$ & $0{,}02$\\
\bottomrule
\end{tabular}}
\end{center}

\begin{center}
\begin{minipage}{0.9\textwidth}
\textbf{Die Logistische zeigt ein Plateau --- $0{,}69 = \ln 2$, acht Schritte lang. Das Rauschen zeigt keines: Die Steigung beginnt riesig und fällt gegen null. Wer bei $k = 1$ abliest, misst $1{,}5$ und nennt lineares Rauschen chaotisch.}
\end{minipage}
\end{center}

\begin{center}
\begin{minipage}{0.88\textwidth}
\kursiv{\textbf{Der Exponent ist eine Rate, und eine Rate braucht ein Plateau.} Rauschen trennt Nachbarn sofort und vollständig --- der Abstand springt im ersten Schritt auf den typischen Abstand zweier beliebiger Punkte und kann dann nicht mehr wachsen: formal $\lambda = \infty$ bei $k \to 0$ und $0$ danach. Ein chaotisches System trennt sie \kursiv{allmählich}, mit konstanter Rate, bis der Attraktor voll ist. Die Diagnose ist also nicht die Zahl, sondern die Gerade: Ohne lineares Stück in der Auftragung gibt es keinen Exponenten --- und die Literatur der neunziger Jahre ist voll von positiven Exponenten aus Herzfrequenz, EEG und Börsenkursen, die bei $k = 1$ abgelesen wurden. Kapitel 12 gibt die Prüfung, die sie bestanden haben müssten.}
\end{minipage}
\end{center}
\end{warnung}

\section{Übungsaufgaben}

\begin{aufgabe}
\boxtitel{Aufgabe 9.1 --- Der Exponent von Hand}
Berechnen Sie den Ljapunow-Exponenten der Zeltabbildung $f(x) = 2x$ für $x < \tfrac12$, $2 - 2x$ sonst, und den der logistischen Abbildung bei $r = 4$ über die Substitution $x = \sin^{2}(\pi\theta/2)$.

\loesung
\schritt{Schritt 1 --- Zelt}
$|f'| = 2$ überall (außer in $\tfrac12$): $\lambda = \lim\frac1n\sum\ln|f'(x_i)| = \ln2$ für jeden Verlauf, der $\tfrac12$ meidet. \checkmark

\schritt{Schritt 2 --- Die Substitution}
$x = \sin^{2}(\pi\theta/2)$, $\theta\in[0,1]$. Dann $4x(1-x) = 4\sin^{2}\cos^{2} = \sin^{2}(\pi\theta)$ $= \sin^{2}(\pi\cdot2\theta/2)$: In $\theta$ ist die Abbildung $\theta\mapsto 2\theta$ für $\theta < \tfrac12$ und, wegen $\sin^{2}(\pi\theta) = \sin^{2}(\pi(1-\theta))$, $\theta\mapsto 2 - 2\theta$ sonst --- das Zelt.

\schritt{Schritt 3 --- Der Exponent überträgt sich}
Die Substitution ist eine glatte Umkehrung mit beschränkter Ableitung (außer an den Rändern); Dehnungsraten im Mittel sind davon unabhängig: $\lambda = \ln 2 = 0{,}693$. \checkmark\ (Tabelle: $0{,}693$.)

\begin{center}
\begin{minipage}{0.88\textwidth}
\kursiv{\textbf{Die Substitution ist von \Person{Ulam} und \Person{von Neumann} 1947} --- sie benutzten $x_{n+1} = 4x_n(1-x_n)$ als Zufallsgenerator (Band 27, Kapitel 2) und wussten, warum er funktioniert: Jeder Schritt schiebt eine Binärstelle von $\theta$ nach vorn. Deterministisch, und doch ein Münzwurf je Schritt --- die Brücke zwischen Band 27 und diesem.}
\end{minipage}
\end{center}
\end{aufgabe}

\begin{aufgabe}
\boxtitel{Aufgabe 9.2 --- Kaplan--Yorke}
Die \Begriff{Ljapunow-Dimension} ist $D_L = j + \frac{\lambda_1 + \cdots + \lambda_j}{|\lambda_{j+1}|}$ mit dem größten $j$, für das die Summe noch $\ge 0$ ist. Berechnen Sie sie für Lorenz und erklären Sie die Formel.

\loesung
\schritt{Schritt 1 --- $j$}
$\lambda_1 = 0{,}906 \ge 0$; $\lambda_1 + \lambda_2 = 0{,}906 \ge 0$; $\lambda_1 + \lambda_2 + \lambda_3 = -13{,}67 < 0$. Also $j = 2$.

\schritt{Schritt 2 --- Zahl}
$D_L = 2 + 0{,}906/14{,}57 = 2{,}062$.

\schritt{Schritt 3 --- Warum}
Ein $j$-dimensionales Volumenelement wächst mit $\mathrm e^{(\lambda_1+\cdots+\lambda_j)t}$. Für $j = 2$ wächst die Fläche (Rate $0{,}906$), für $j = 3$ schrumpft das Volumen (Rate $-13{,}67$): Der Attraktor ist \enquote{mehr als eine Fläche, weniger als ein Körper}, und die Interpolation $j + (\text{Rest})/|\lambda_{j+1}|$ ist die Dimension, bei der Wachsen und Schrumpfen sich die Waage halten. \Person{Kaplan} und \Person{Yorke} vermuteten 1979, dass das die Dimension des Attraktors ist; für Lorenz liefert die direkte Messung (Kapitel 11) $2{,}05 \pm 0{,}01$. \checkmark
\end{aufgabe}

\begin{aufgabe}
\boxtitel{Aufgabe 9.3 --- Wie viele Stellen für ein Jahr}
Ein Fluss hat $\lambda_1 = 0{,}35$ je Tag. Die Vorhersage soll ein Jahr reichen bei $\Delta/\delta_0$ als Verhältnis von erlaubtem Fehler zu Anfangsfehler. Welches Verhältnis braucht man, und ist das möglich?

\loesung
\schritt{Schritt 1}
$\ln(\Delta/\delta_0) = \lambda_1 t = 0{,}35\cdot365 = 127{,}8$.

\schritt{Schritt 2}
$\Delta/\delta_0 = \mathrm e^{127{,}8} = 10^{55{,}5}$: der Anfang müsste auf $55$ Dezimalstellen genauer bekannt sein als der erlaubte Fehler am Ende.

\schritt{Schritt 3 --- Einordnung}
Das Verhältnis von Erddurchmesser zu Protondurchmesser ist $10^{22}$; $10^{55}$ liegt jenseits jeder physikalischen Messung, und jenseits der Gültigkeit der Gleichungen, die bei solchen Skalen durch andere ersetzt werden müssten.

\begin{center}
\begin{minipage}{0.88\textwidth}
\kursiv{\textbf{Das ist der Sinn von \enquote{prinzipiell unvorhersagbar}}: nicht, dass man es nicht weiß, sondern dass die Zahl der nötigen Stellen jede Physik übersteigt. Es ist auch der Grund, weshalb Wettervorhersage als Ensemble betrieben wird (Kapitel 15): Nicht einen Anfang vorhersagen, sondern fünfzig leicht verschiedene, und die Streuung der Ergebnisse als das nehmen, was sie ist --- die Form des Nichtwissens nach $t$ Tagen.}
\end{minipage}
\end{center}
\end{aufgabe}

\begin{geschichte}
\boxtitelg{Ljapunows Dissertation und Oseledets' Satz}
\Person{Alexander Ljapunow} reichte 1892 in Charkiw seine Dissertation \enquote{Das allgemeine Problem der Stabilität der Bewegung} ein; die Exponenten, die seinen Namen tragen, dienten dort als Werkzeug, um Stabilität zu entscheiden, ohne die Lösungen zu kennen --- das Gegenstück zur Funktion aus Satz 2.4.

\begin{center}
\begin{minipage}{0.9\textwidth}
\kursiv{Dass sie für chaotische Systeme eine Rate des Vergessens messen, war nicht sein Thema; er wollte zeigen, dass Bewegungen stabil sind. \textbf{Ihre Existenz für fast alle Anfangswerte bewies \Person{Oseledets} 1965 in Moskau} --- mit einem Satz, dessen Beweis so dicht war, dass \Person{Raghunathan} 1979 und \Person{Ruelle} 1979 ihn neu führten, um ihn benutzbar zu machen. \Person{Benettin}, \Person{Galgani}, \Person{Giorgilli} und \Person{Strelcyn} gaben 1980 das Verfahren aus Satz 9.3, mit dem seither jeder Exponent gerechnet wird; \Person{Wolf} und Mitarbeiter 1985 das Gegenstück für Daten.}
\end{minipage}
\end{center}

\schritt{Zwei Wochen}
\begin{center}
\begin{minipage}{0.9\textwidth}
\kursiv{Die Grenze von zwei Wochen für das Wetter hat \Person{Lorenz} 1969 aus der Fehlerverdopplung in einem Modell mit vielen Skalen abgeleitet: Kleine Wirbel vergessen in Stunden, große in Tagen, und der Fehler wandert von klein nach groß. \textbf{Fünfzig Jahre später, 2019, haben \Person{Zhang}, \Person{Emanuel} und Mitarbeiter mit den besten Modellen dieselben zwei Wochen gefunden}, mit der Aussage, dass selbst eine perfekte Beobachtung im Kleinen die Grenze nur um wenige Tage verschiebt. Lorenz' Schätzung mit Bleistift und Papier hat gehalten.}
\end{minipage}
\end{center}
\end{geschichte}

\begin{forschung}
\boxtitelp{Exponenten heute}
\begin{center}
{\small
\renewcommand{\arraystretch}{1.45}
\begin{tabular}{@{}P{3.0cm}P{5.2cm}@{}}
\toprule
\textbf{Kovariante Ljapunow-Vektoren} & Die Richtungen zu den Exponenten --- lokal, entlang des Verlaufs (Ginelli 2007): Wo im Zustandsraum ist das System am empfindlichsten?\\
\textbf{Lokale Exponenten} & $\lambda$ über endliche Zeit, als Funktion des Ortes: Vorhersagbarkeit von Fall zu Fall --- die Grundlage der Ensemblevorhersage.\\
\textbf{Große Systeme} & Spektrum von Tausenden Exponenten in Turbulenz und neuronalen Netzen; die Zahl der positiven als Maß der Freiheitsgrade des Chaos.\\
\bottomrule
\end{tabular}}
\end{center}

\schritt{Zur zweiten Zeile}
\begin{center}
\begin{minipage}{0.88\textwidth}
\kursiv{\textbf{Der Exponent aus Definition 9.1 ist ein Mittel über den ganzen Attraktor, und das Mittel verbirgt, dass Chaos nicht überall gleich stark ist.} Auf dem Lorenz-Attraktor ist die Trennung nahe dem Ursprung heftig und auf den Flügeln sanft; eine Vorhersage, die am Flügel beginnt, hält länger als eine, die am Wechsel beginnt. Die Wettervorhersage nutzt das seit den neunziger Jahren: Sie rechnet nicht nur den Zustand, sondern auch, wie empfindlich er gerade ist --- und sagt an manchen Tagen eine Woche, an anderen drei Tage voraus. Der Vorhersagehorizont ist selbst vorhersagbar.}
\end{minipage}
\end{center}
\end{forschung}

\begin{merksatz}
\boxtitel{Was Sie jetzt können}
\begin{center}
\renewcommand{\arraystretch}{1.5}
\begin{tabular}{@{}P{3.0cm}P{5.6cm}@{}}
\toprule
\textbf{Definition 9.1} & Ljapunow-Exponenten\\
\textbf{Satz 9.2} & Vorzeichen: Fixpunkt, Zyklus, Torus, Chaos; Summe = Divergenz\\
\textbf{Satz 9.3} & Benettin: normieren und mitteln\\
\textbf{Satz 9.4} & Horizont $= \ln(\Delta/\delta_0)/\lambda_1$\\
\textbf{Der Bruchfall} & ohne Plateau kein Exponent; Rauschen hat keines\\
\bottomrule
\end{tabular}
\end{center}

\begin{center}
\begin{minipage}{0.88\textwidth}
\kursiv{Teil III ist zu Ende: Die Verdopplung führt ins Chaos, drei Gleichungen reichen, und der Exponent sagt, wie weit man sieht. Teil IV kehrt die Frage um: Man hat keine Gleichungen, nur Daten --- eine Herzfrequenz, ein EEG, einen Kurs. Kann man ihnen ansehen, ob dahinter ein Attraktor steckt, und wie groß er ist?}
\end{minipage}
\end{center}
\end{merksatz}
